EGL and HGL

HGL (hydraulic grade line):
$$HGL=z+\frac{p}{\rho g}$$
EGL (energy grade line):
$$EGL=HGL+\frac{\alpha_E\bar V^2}{2g}$$

Vertical separation EGL-HGL = velocity head.
Pressure head = HGL - pipe centerline elevation.

NO MACHINES, IDEAL:
EGL constant, despite bends, rise, or contraction.

VISCOUS FLOW:
EGL drops along pipes and across fitting losses.
PUMP: EGL jumps UP.
TURBINE: EGL drops DOWN.

LARGE TANK FREE SURFACE:
Gauge p=0, V about 0, so EGL and HGL meet at free-surface elevation.

ATMOSPHERIC FREE JET:
HGL = jet z; EGL stays above by velocity head.

CONTRACTING/RISING PIPE SKETCH
1. Draw pipe centerline z(s).
2. Continuity: V increases when diameter shrinks.
3. Sketch EGL from energy and losses.
4. Construct HGL=EGL-velocity head.
5. Static pressure is rho*g*(HGL-z).

In a contraction, EGL-HGL gap enlarges.
A sudden expansion may make HGL rise even though EGL falls.
HGL below centerline means negative GAUGE p, not automatically vaporization.
